\begin{proof}[Proof of \cref{obj:thm:nonlocal-testing-answer}]
\begin{enumerate}
\item The canonical design is the product of the treatment-assignment and audit-mark laws. On the finite population \(V\), with \(|V|=n\), its marginals therefore satisfy
\[
\operatorname{Law}(Z)
=\bigotimes_{j\in V}\operatorname{Bernoulli}(1/2),
\qquad
\operatorname{Law}(W)
=\bigotimes_{\substack{(j,i)\in V^2\\j\ne i}}
\operatorname{Bernoulli}(q).
\]
Indeed, \(q\in[0,1]\) ensures that both factors are probability laws, and projecting their product onto either coordinate gives the corresponding factor. Thus the design satisfies \cref{obj:ass:assignment,obj:ass:audit}.
% lean: CausalSmith.Experimentation.ThinnedgraphAdditiveRiskFrontier.thinnedDesign_assignment
% lean: CausalSmith.Experimentation.ThinnedgraphAdditiveRiskFrontier.thinnedDesign_audit
The two coordinate projections of a product probability law are independent, so
\[
W\mathbin{\perp\!\!\!\perp}Z.
\]
This verifies \cref{obj:ass:design-independence}.
% lean: CausalSmith.Experimentation.ThinnedgraphAdditiveRiskFrontier.thinnedDesign_independent

\item Apply \cref{obj:thm:block-testing-scale} to this design. Its population, block, degree, capacity, and retention hypotheses are precisely the hypotheses here, and the preceding step verifies its three design hypotheses. Its complete-record mixture laws, testing scale, and testing radius coincide with \(P_{+1,h},P_{-1,h},s,\bar h\) in the present statement. Consequently, it supplies
\[
(100\pi)^{-1}s\le\bar h,
\qquad
\bar h\le2s,
\]
as well as
\[
\forall h\in\mathbb R,\qquad
0\le h\le(100\pi)^{-1}s
\ \Longrightarrow\
\operatorname{TV}(P_{+1,h},P_{-1,h})\le\frac12.
\]
% lean: CausalSmith.Experimentation.ThinnedgraphAdditiveRiskFrontier.block_testing_scale

\item Suppose \(d=2\). Expanding the definition of the association-revelation probability gives
\[
m=2B,
\qquad
p=1-(1-q)^2=2q-q^2.
\]
First suppose \(q>0\). Since \(q\le1\) and \(B\ge1\),
\[
p=q(2-q)>0,
\qquad
2Bp>0,
\qquad
\sqrt{2Bp}>0.
\]
The definition of the testing scale therefore gives
\[
s=\min\left\{1,\frac{2}{\sqrt{2Bp}}\right\}.
\]
To simplify the second entry, square it:
\[
\left(\frac{2}{\sqrt{2Bp}}\right)^2
=\frac{4}{2Bp}
=\frac{2}{Bp}.
\]
Both the quotient and the square root below are nonnegative, so
\[
\frac{2}{\sqrt{2Bp}}
=\sqrt{\frac{2}{Bp}},
\qquad
s=\min\left\{1,\sqrt{\frac{2}{B(2q-q^2)}}\right\}.
\]
In the remaining case \(q\le0\), the hypothesis \(q\ge0\) gives
\[
q=0,\qquad p=0,\qquad s=1.
\]
These two cases establish the stated degree-two formula.
% lean: CausalSmith.Experimentation.ThinnedgraphAdditiveRiskFrontier.nonlocal_testing_answer
\end{enumerate}
\end{proof}
